\section{Datos sintéticos y Action Calling: construcción de un volante de entrenamiento escalable}

\subsection{De las tareas de percepción a las tareas de acción}
El expositor enfatiza que la comprensión visual por sí sola (entender la interfaz) no basta para formar un Agente robusto; la capacidad debe transferirse a «invocar acciones». Por ello, el curso introduce un pipeline de datos sintéticos que convierte preguntas y respuestas sobre imágenes, OCR, localización y razonamiento de flujos en muestras ejecutables como acciones.

\begin{importantbox}{La esencia de Action Calling}
Action Calling no es una función adicional, sino el cambio del espacio de salida del modelo de «describir» a «ejecutar». Esto exige que los datos de entrenamiento incluyan a la vez intención, ruta de razonamiento, formato de la acción y retroalimentación de la ejecución.
\end{importantbox}

\begin{figure}[H]
\centering
\includegraphics[width=0.88\textwidth]{frame-06-synthetic-data.jpg}
\caption{Video aproximadamente en 01:02:32: el pipeline de datos sintéticos unifica preguntas visuales, cadenas de razonamiento e invocación de acciones en un mismo flujo de entrenamiento}
\end{figure}

\subsection{Flujo sintético típico}
\begin{enumerate}
    \item Muestrear escenarios de tarea y estados de la interfaz;
    \item Construir preguntas y objetivos (que incluyan localización, lectura, operación y verificación);
    \item Usar el modelo para generar candidatos de razonamiento y secuencias de acciones;
    \item Filtrar muestras de alta calidad con ayuda del ejecutor y del evaluador;
    \item Pasar a un entrenamiento por etapas y retroalimentar continuamente las muestras fallidas.
\end{enumerate}

\begin{knowledgebox}{Por qué funciona el entrenamiento por etapas}
En el curso se menciona varias veces: «primero aprender a ver, luego aprender a hacer, y después aprender la estrategia de largo plazo». Si desde el inicio se aprende de extremo a extremo una acción de cadena larga, el modelo cae con mayor facilidad en alta varianza y óptimos locales. Dividir el entrenamiento en etapas reduce la dificultad de la optimización y mejora el aprovechamiento de las muestras.
\end{knowledgebox}

\subsection{Los límites de los datos sintéticos}
Aunque las muestras sintéticas permiten escalar, el sesgo de distribución es inevitable. El curso recomienda mantener equilibrada la proporción entre trayectorias reales y trayectorias sintéticas, y hacer de manera continua pruebas de regresión con tareas en producción, para evitar que el modelo solo se desempeñe bien dentro de la distribución sintética.

\begin{warningbox}{Los tres grandes riesgos de los datos sintéticos}
\begin{itemize}
    \item Desvío semántico: la descripción del objetivo no coincide con la intención real de la tarea;
    \item Alucinación de acciones: se generan operaciones que parecen razonables pero no son ejecutables;
    \item Acumulación de sesgo: los datos generados por el modelo terminan reforzando sus propios patrones de error.
\end{itemize}
\end{warningbox}

\subsection{Resumen de la sección}
El uso correcto de los datos sintéticos es como «complemento y aceleración», no como sustituto de los datos reales de interacción. La experiencia clave que ofrece el curso es: convertir la verificación de la ejecución en una restricción dura del pipeline sintético.
